\subsection*{README Generation}

To enhance code readability and accessibility, README files were written for the project.
The README writing was partly assisted by AI.

\textbf{Prompt:}
\begin{lstlisting}[language={}, breaklines=true]
I am developing an embedded 2D game on an ESP32-S3. The application has two main components with vastly different timing requirements: a high-frequency game loop running at circa 50Hz (handling low-pass filtered IMU input, physics sub-stepping, collision resolution, and display rendering) and a low-frequency network stack running at circa 2Hz (handling blocking MQTT connections, TLS handshakes, and telemetry JSON serialization).

Use the provided context to append to, enhance and update the README.md file for the project with the most recent changes.
Use headings, subheadings, bullet points, inline code, notes, and other formatting to make the README.md file more readable and keeping it professional.
Do not change the provided README structure, but enhance it with the most recent changes and improvements.

Summarize the code changes and provide a brief overview about the usage.
\end{lstlisting}

\textbf{Context:}
\begin{lstlisting}[language={}, breaklines=true]
Current file (cpp/h), README.md text snippets, and any other relevant information can be used to generate the docstring.
\end{lstlisting}

\textbf{Answer:}
\begin{lstlisting}[language={}, breaklines=true]
Generated README.md content with the most recent changes and improvements.
\end{lstlisting}
